\section{The selected first-order arithmetic bound}
\label{sec:count-first-order-v037}

\status{CFO1--CFO21 supply the original first-order bound
\texttt{B1} on the same full-block, $g=0$ \texttt{Kap} produced by
CCH15.  For the actual rank-four Count-II application the exponent is
exactly $4[L^+:\mathbb Q]+4q$.  The proof includes the uniform dual
cocycle estimate, the original arithmetic presentation at every
Taylor--Wiles level, and the complete local/global cardinality
cancellation.  The original \texttt{CohP}, \texttt{LocCFT}, old
$D_0$ condition and arithmetic data are retained.  This closes the
selected first-order supplier; it does not assert the universal
arbitrary-\texttt{Kap} predicate or supply the remaining ordinary
automorphic and same-ideal support producers.}

We use the pinned source conventions throughout.  For an additive
finite module $M$ with a continuous action $\alpha$ of a profinite
group $\Gamma$, put
\[
 \begin{split}
 Z^1(\Gamma,M)&=\{z:\Gamma\longrightarrow M:
 z\text{ is locally constant},\
 z(gh)=z(g)+\alpha(g)z(h)\},\\
 B^1(\Gamma,M)&=\{\delta Y:g\longmapsto\alpha(g)Y-Y\}.
 \end{split}                                                    \tag{CFO.1}
\]
These are the literal source \texttt{Zl} and \texttt{Bl}, after
forgetting the redundant module structure.  All cardinalities below
refer to finite sets; their finiteness is established before division.

% BEGIN CFO1
\paragraph{CFO1. The selected objects and the exact target.}
\label{claim:CFO1}
Fix an odd prime $\ell$, the original CM field $L$ and algebraic
closure $\Omega$, $F=L^+$, an original
$s:\texttt{HidaSetting}(\ell,L,\Omega)$, $V_0:s.\texttt{Lev}$ with
$V_0.S'=s.S$, and the original coefficient DVR $O$ and
$r:\texttt{Resid}(V_0.\texttt{frame},O,n)$.  Here $n=4$ in both
Count-II links.  Let $T:\texttt{TWGrp}(\texttt{globOf}(r,V_0))$ and
$P:\texttt{TWPl}(r,T)$ be the actual output of CAP13 for the finite
list chosen by CCH15.  Let $x$ and $\kappa$ be that same extended
\texttt{NP} and its chosen \texttt{Kap}.  Thus $x.a=n$,
$\kappa.g=0$, and the roots $\alpha_{ji}$ give the literal
$T.\texttt{Eig}(x,\alpha)$.  Write
\[
 \begin{gathered}
 \Gamma_0=V_0.\texttt{frame}.\Gamma,\quad
 G_0=\ker\eta,\quad \Gamma_N=(T.GN(N)).\Gamma,\quad
 p_N=(T.h(N)).p,\quad G_N=\ker\eta_N,\\
 k=x.K,\quad A=k[[t]],\quad E=k((t)),\quad A_m=A/(t^m),
 \quad Q=|k|,\quad q=T.q,\quad d=[F:\mathbb Q],\\
 r_g=\operatorname{mat}(x.\rho(g)),\quad
 \chi=\epsilon\bmod\ell\text{, embedded as constants in }A,
 \quad \sigma(g)=cgc^{-1},\quad s_c=-\chi(c).
 \end{gathered}                                                   \tag{CFO.2}
\]
The symbol $r_g$ is unrelated to the name of the residual datum $r$.
By CCH1--CCH3, the literal first-order space is
$M_m=\operatorname{Mat}_n(A_m)$ and its paired dual is
$M'_m=\eta\operatorname{Mat}_n(A_m)$; we identify the latter with
its matrix of $\eta$ coefficients.  The matrix $\rho_1$ is $r_g$
in constant dual-number coordinates.  In particular $g=0$ does
not make $M_m$, \texttt{Sel1}, or its deformation classes zero.

Retain the complete original \texttt{CohP} and \texttt{LocCFT}(s),
the residual cyclotomic nontriviality hypothesis, and
\[
 \begin{gathered}
 r_g=1\quad(g\in G_0,\ g\in\operatorname{PL}_\ell(T.H)),\qquad
 x.N4(D_0),\qquad D_0=V_0.\texttt{frame}.D0,\\
 \operatorname{Detects}\bigl(\kappa,\operatorname{range}(T.sg)\bigr).
 \end{gathered}                                                   \tag{CFO.3}
\]
The first line is supplied, where needed in the consumer, by the
original residual-determinant argument recorded in JSP11 and
CCH17; its hypothesis is not replaced by a statement about a
different subgroup.  In particular $D_0$ here is the old
distinguished decomposition image, although the intermediate
\texttt{globOf} frame records a bottom $D0$.

The target is exactly
\[
 \exists C\in\mathbb N\ \forall m\in\mathbb N\
 \forall^\infty N:\quad
 \begin{cases}
  \operatorname{Sel1}_N(A_m,\kappa)\text{ is finite},\\
  |\operatorname{Sel1}_N(A_m,\kappa)|
   \le |\operatorname{Bd1}_N(A_m,\kappa)|\,Q^{m(nd+nq)+C}.
 \end{cases}                                                     \tag{CFO.4}
\]
Here $C$ may depend on the fixed point, list and arithmetic tuple;
it does not depend on $m$ or $N$.  The eventual bound on $N$ may
depend on $m$.  This is \texttt{B1}, with its original quantifier
order, on the chosen \texttt{Kap}.  CCH13 preserves (CFO.3) under
the same finite coefficient extension, and CCH16 already supplies
\texttt{B2} on this identical \texttt{Kap}.
% END CFO1

% BEGIN CFO2
\paragraph{CFO2. Finite modular cocycles and the complete function module.}
\label{claim:CFO2}
Assume the original $\Phi_\ell(\Gamma_0)$, supplied from
\texttt{CohP} by CAP1--CAP2.  If $U$ is an open subgroup of
$\Gamma_0$ and $V$ a finite continuous $\mathbb F_\ell[U]$-module,
then $Z^1(U,V)$ is finite.  Indeed the kernel $U_1$ of its finite
action is open.  Restriction of a cocycle to $U_1$ is an additive
homomorphism into $V$.  Choose an $\mathbb F_\ell$ basis of $V$.
Its coordinate homomorphisms have open kernels, and
$\Phi_\ell(\Gamma_0)$ gives only finitely many possibilities for
each, because $U_1$ is still open in $\Gamma_0$.  A cocycle is
determined by this restriction and its values on a fixed finite
set of representatives for $U/U_1$.  Both choices are finite.
This proves the assertion without a finite-generation assumption
on $U$.

Let $\mathcal C$ be the $A$-module of functions
$f:G_0\longrightarrow\operatorname{Mat}_n(A)$ whose reduction
modulo every $t^a$ is locally constant.  It is $t$-torsion-free,
separated and complete.  For completeness, a $t$-adically Cauchy
sequence has, modulo each $t^a$, an eventually constant locally
constant function.  These compatible functions lift entrywise to
one $A$-valued function.  Moreover
\[
 f(G_0)\subset t^a\operatorname{Mat}_n(A)
 \quad\Longleftrightarrow\quad f\in t^a\mathcal C:                 \tag{CFO.5}
\]
division by $t^a$ preserves the required local constancy by
looking at reduction modulo $t^{a+b}$.  Every locally constant
function with values in $\operatorname{Mat}_n(A_m)$ lifts to
$\mathcal C$: choose the finitely many polynomial representatives
of its finitely many values.

The full twisted cocycles at level $m$ satisfy
\[
 z(gh)=z(g)+\chi(g)r_gz(h)r_g^{-1},\qquad
 z(\sigma(g))=s_c z(g)^t.                                        \tag{CFO.6}
\]
They are a subset of the finite cocycles just counted, for the
finite twisted matrix action of $G_0$.  This proves their
finiteness for every $m$, including the zero module at $m=0$.
It also proves finiteness of every subset obtained by imposing
the finite trace equations below.  These are precisely the
modular \texttt{TwZ} equations of CCH3.
% END CFO2

% BEGIN CFO3
\paragraph{CFO3. A complete-DVR approximation and counting lemma.}
\label{claim:CFO3}
Let $\mathcal C$ be a complete, separated, $t$-torsion-free
$A$-module and $\mathcal C'$ a separated, $t$-torsion-free
$A$-module.  Let $\Phi:\mathcal C\to\mathcal C'$ be linear, set
$K=\ker\Phi$, and let $B\subset K$ be a submodule.  Suppose
\[
 S_1/t\mathcal C\text{ is finite},\quad
 S_a=\Phi^{-1}(t^a\mathcal C'),\quad
 \forall z\in K\ \exists e:\ t^e z\in B.                          \tag{CFO.7}
\]
Then there is $b_0\in\mathbb N$ such that, for every $a\ge0$,
\[
 |S_a/t^a\mathcal C|
 \le |(B+t^a\mathcal C)/t^a\mathcal C|\,Q^{b_0}.                  \tag{CFO.8}
\]
All the quotients occurring here are finite.

First, $S_a/t^a\mathcal C$ is finite by induction.  Its reduction
modulo $t$ has finite image in $S_1/t\mathcal C$; after choosing
one representative in a fiber, divide differences by $t$.
Torsion-freeness of $\mathcal C'$ embeds that fiber into
$S_{a-1}/t^{a-1}\mathcal C$.  The case $a=0$ is a singleton.
The images of $S_a$ in $\mathcal C/t\mathcal C$ therefore form a
descending sequence of finite sets, stationary from some $b\ge1$.

We spell out the approximation consequence of this stationarity.
For every $r\ge1$, $x\in S_{b+r-1}$ and integer $M$ there is
$y\in S_M$ with $x-y\in t^r\mathcal C$.  For $r=1$ this is
stationarity (or use $x$ when $M\le b$).  Suppose it holds at $r$.
For $x\in S_{b+r}$ choose, by the $r=1$ assertion,
$z\in S_L$ with $x-z=t u$, where
$L\ge\max(M,b+r)$.  Division in $\mathcal C'$ gives
$u\in S_{b+r-1}$.  Apply the induction hypothesis to find
$v\in S_M$ with $u-v\in t^r\mathcal C$.
Then $y=z+tv\in S_M$ and $x-y\in t^{r+1}\mathcal C$.

Starting with $x_0\in S_{b+j}$, apply this consequence successively
to obtain
\[
 x_i\in S_{b+j+i},\qquad
 x_{i+1}-x_i\in t^{j+i+1}\mathcal C.
 \]
Completeness gives a limit $k_0\in\mathcal C$.  For every $a$,
linearity and a sufficiently large $i$ imply
$\Phi(k_0)\in t^a\mathcal C'$, so separation gives $k_0\in K$.
Since $x_0-k_0\in t^j\mathcal C$, division by $t^j$ also gives
\[
 S_{b+j}\subset K+t^j S_b.                                      \tag{CFO.9}
\]

The same division argument gives
$K\cap t^a\mathcal C=t^aK$.  Thus $K/tK$ injects into the finite
$S_1/t\mathcal C$.  Choose finitely many lifts $k_1,\ldots,k_h$
of generators of $K/tK$.  Repeatedly subtract their residues and
divide by $t$.  For any $k\in K$ this gives
$k=\sum_i a_i k_i$ with $a_i\in A$: the coefficients converge
in $A$ and the residual term lies in every $t^a\mathcal C$.
Equivalently, $K$ is a closed complete submodule with this
explicit topological Nakayama proof.  It is finitely generated
as an ordinary $A$-module.  By (CFO.7) a single $e$ has
$t^eK\subset B$.  Hence $K/B$ is finite.

Choose finite representatives $\mathcal G$ of $K/B$ and
$\mathcal F$ of $S_b/t^b\mathcal C$.  For $a=b+j$, (CFO.9)
shows that the classes of
$g+t^j f$, $g\in\mathcal G$, $f\in\mathcal F$, cover
$S_a/(B+t^a\mathcal C)$.  This quotient therefore has at most
$|\mathcal G||\mathcal F|$ elements.  The finitely many cases
$a<b$ have finite quotients too.  Choose $b_0$ large enough to
bound all these integers by $Q^{b_0}$; $Q\ge2$.  Taking the
finite fibers of the quotient by the image of $B$ proves
(CFO.8).  No bound has been chosen after fixing $a$.
% END CFO3

% BEGIN CFO4
\paragraph{CFO4. The uniform annihilated-cocycle estimate.}
\label{claim:CFO4}
Use the complete module $\mathcal C$ of CFO2.  Define a linear
map $\Phi$ into the product of copies of $\operatorname{Mat}_n(A)$
and $A$ whose coordinates are the defects in (CFO.6), for all
$g,h$, and
\[
 \operatorname{tr}\bigl(f(s_j)r_{s_j}\bigr),\qquad s_j=T.sg(j).
                                                                    \tag{CFO.10}
\]
The target product is separated and $t$-torsion-free.  By (CFO.5)
and the finite-value lifting assertion, $S_m/t^m\mathcal C$ is
exactly the set $\mathcal A_m$ of modular \texttt{TwZ} annihilated
by all the traces (CFO.10).  It is finite by CFO2, so the first
hypothesis of CFO3 holds.

Let $B$ consist of the integral functions
\[
 g\longmapsto\chi(g)r_gYr_g^{-1}-Y,\qquad s_cY^t=Y.               \tag{CFO.11}
\]
These are the literal integral \texttt{TwB}, with its transpose
condition.  They lie in $\ker\Phi$: the cocycle and polarization
identities follow from the original form, and $\chi(s_j)=1$
because each $s_j$ is in the full cyclotomic kernel.  Cyclicity
of matrix trace kills (CFO.10).  The original \texttt{Detects}
condition, supplied on the finite list by CCH15, is exactly
the final hypothesis in (CFO.7).  Hence CFO3 gives one $C_2$
such that
\[
 |\mathcal A_m|\le|\operatorname{TwB}_m|Q^{C_2}
 \qquad(m\ge0).                                                  \tag{CFO.12}
\]
Here the image of the integral $B$ is contained in the modular
\texttt{TwB}; equality is unnecessary.  In particular the proof
does not assume that a modular solution lifts to an integral
solution of the equations.  CFO3 is precisely the uniform
approximation argument that supplies this otherwise missing
step.  This reconstructs the source annihilated-cocycle
cardinality endpoint at the selected scope.
% END CFO4

% BEGIN CFO5
\paragraph{CFO5. Pro-$\ell$ kernels and exact descent from a level.}
\label{claim:CFO5}
For a subgroup $H$ of a compact group, the literal
$\operatorname{PL}_\ell(H)$ is the intersection of all open
normal $U$ such that each $h\in H$ has an $\ell$-power in $U$.
These $U$ form a downward directed family: for two of them,
add the two exponents for each $h$.  Compactness implies that
any open set containing this intersection contains one such
$U$; otherwise their closed intersections with its compact
complement have the finite intersection property.  Consequently
each $h\in H$ has an $\ell$-power in that open set.

Apply this to the open kernel of $r\bmod t^m$ in $G_0$.
It is also an open subset of $\Gamma_0$, and contains
$\operatorname{PL}_\ell(T.H)$ by (CFO.3).
Here $\operatorname{PL}_\ell(T.H)\subset G_0$ follows
directly by using $G_0$ itself among the open normal
subgroups in its defining intersection.
It follows that $r_h\bmod t^m$ has $\ell$-power order for
$h\in T.H$.  Since $\chi(h)=1$, the twisted adjoint action
of such $h$ on $M'_m$ has $\ell$-power order too.

For a finite module action $\beta$ and a continuous cocycle $y$,
the maps
$v\mapsto y(g)+\beta(g)v$ form a continuous homomorphism to the
finite permutation group of the module.  Its kernel is an open
normal subgroup, given exactly by
$\beta(g)=1$ and $y(g)=0$.  If $\beta(h)$ has $\ell$-power
order, then the affine permutation has $\ell$-power order:
after that power it is a translation, whose $\ell$th power is
the identity.  Thus the affine kernel contains
$\operatorname{PL}_\ell(H)$.

Now take any action $\beta$ on $M'_m$ satisfying the source
\texttt{IsAd} on $\Gamma_N$, and a cocycle $y$ with
$y(c_N)=0$ and $y(T.ta(N,j))=0$ for every $j$.  Its affine
action is continuous: the \texttt{IsAd} formulas and the
finite reduction of $r$ determine a locally constant
action on each of the two open cosets.  Its affine
kernel contains the level
$\operatorname{PL}_\ell(p_N^{-1}T.H)$ by the preceding argument.
It contains every $T.ta(N,j)$ as well, because $p_N(ta)=1$
makes its linear action the identity.  Being closed and normal,
it contains the closure of their normal closure together with
that pro-$\ell$ kernel.  The original \texttt{TWGrp.ker\_le}
therefore places $\ker p_N$ inside the affine kernel.

The cocycle factors through $p_N$.  Its descended function is
locally constant: a continuous surjection from a compact group
to a Hausdorff group is a quotient map, and the fibers of the
descended function have open preimages.  Restriction to $G_0$
gives a modular \texttt{TwZ}.  Indeed the cocycle equation gives
the first identity in (CFO.6), and normalization at $c_N$ gives
$z(\sigma g)=\beta(c)z(g)=s_cz(g)^t$.  Conversely this restriction
determines the normalized cocycle, since every element of
$\Gamma_N$ is uniquely $g c_N^e$, with $g\in G_N$ and $e=0,1$.
This proves the complete selected descent statement, including
continuity and its injectivity assertion.
% END CFO5

% BEGIN CFO6
\paragraph{CFO6. Eventual Frobenius comparison for every modular cocycle.}
\label{claim:CFO6}
Fix $m$.  For any modular $z$ satisfying (CFO.6), the affine
action on $G_0$ has an open normal kernel $U_z$.  We need an
open normal subgroup of the full $\Gamma_0$ containing
$\operatorname{PL}_\ell(T.H)$, not just a subgroup of $G_0$.
Here are the details.

For $h\in T.H$, CFO5 gives an $\ell$-power in the kernel of
the linear action, and one extra $\ell$-power kills the affine
translation.  Thus $h$ has an $\ell$-power in $U_z$.  Take the
normal core of $U_z$ in $\Gamma_0$.  There are only finitely
many conjugates since $U_z$ is open.  Because the original
$T.H$ is normal, the same power assertion holds in each
conjugate, and then in their intersection.  By the definition
of $\operatorname{PL}_\ell$, this core contains that subgroup.
Intersect also with the cores of the finite linear-action
kernel and $\ker\chi$ if necessary.

The original \texttt{ph\_lim} now implies, for this $z$ and
every $j$, eventual equality of the finite linear action,
cyclotomic value and cocycle value at $p_N(ph_{Nj})$ and $s_j$.
For the cocycle value, write $p_N(ph_{Nj})=u s_j$ with $u$
in the affine kernel and use the cocycle equation.  There are
finitely many $j$ and, by CFO2, finitely many modular
\texttt{TwZ}.  Intersect these finitely many eventual sets.
Also apply the same argument to the finite matrix
representation $r\bmod t^m$ itself.  We obtain one eventual
set of $N$, depending on $m$ but on no later cocycle or action,
such that
\[
 \begin{gathered}
 r_{p_N(ph_{Nj})}=r_{s_j}\quad\text{in }\operatorname{Mat}_n(A_m),
 \qquad \chi(p_N(ph_{Nj}))=1,\\
 z(p_N(ph_{Nj}))=z(s_j)
 \quad\text{for every modular \texttt{TwZ} }z\text{ and every }j .
 \end{gathered}                                                   \tag{CFO.13}
\]
This is the full selected \texttt{phEq} statement.  It uses
the literal approximation in \texttt{TWGrp}; no replacement of
inverse Frobenius by arithmetic Frobenius has been made.
% END CFO6

% BEGIN CFO7
\paragraph{CFO7. The uniform dualN estimate, with all its quantifiers.}
\label{claim:CFO7}
There exists the single $C_2$ of CFO4 such that, for every $m$
and eventually every $N$, the following holds for every
representation $\beta$ satisfying the exact \texttt{IsAd}
condition on $M'_m$, provided $Z^1(\Gamma_N,M'_m)$ is finite.
Let
\[
 \mathcal S_\beta=
 \left\{y\in Z^1(\Gamma_N,M'_m):
 \ \forall j\ \exists Y_j\
 \begin{matrix}
 y(ph_{Nj})=\beta(ph_{Nj})Y_j-Y_j,\\
 y(ta_{Nj})=\beta(ta_{Nj})Y_j-Y_j
 \end{matrix}\right\}.
 \]
The same $Y_j$ is used in the two equations.  Then
\[
 |\mathcal S_\beta|\le |B^1(\Gamma_N,M'_m)|Q^{C_2}.                \tag{CFO.14}
\]

The set $\mathcal S_\beta$ is an additive subgroup containing
all coboundaries.  The action on $ta_{Nj}$ is the identity, so
all its elements vanish there.  On the subgroup
$\mathcal S_{\beta,c}$ of those with $y(c_N)=0$, CFO5 and
CFO6 give an injection into $\mathcal A_m$.  To check the trace
equation, the value at $s_j$ is
$r_{s_j}Y_jr_{s_j}^{-1}-Y_j$ by (CFO.13); multiplication by
$r_{s_j}$ and cyclicity of trace make (CFO.10) zero.
Every modular twisted coboundary (CFO.11) is the restriction of
a normalized global coboundary, by inflating its same $Y$.
Consequently
\[
 |\mathcal S_{\beta,c}|
 \le |\mathcal A_m|
 \le |\operatorname{TwB}_m|Q^{C_2}
 \le |B^1(\Gamma_N,M'_m)_c|Q^{C_2}.                              \tag{CFO.15}
\]

For clarity, the normalization factors cancel exactly.
For a cocycle $y$, $(1+\beta(c_N))y(c_N)=0$.
Because $2$ is invertible, this anti-invariant subspace is
$(\beta(c_N)-1)M'_m$.  Evaluation at $c_N$ from any additive
subgroup containing all coboundaries has precisely that image.
Thus $\mathcal S_\beta$ and $B^1$ have the same evaluation
image and
\[
 \frac{|\mathcal S_\beta|}{|B^1|}
   =\frac{|\mathcal S_{\beta,c}|}{|B^1_c|}.
 \]
The denominators are positive because they contain zero.
This proves (CFO.14).  The two \texttt{IsAd} formulas on
$G_N$ and $c_N$ determine the whole action; the eventual
set in CFO6 is therefore independent of the chosen presentation
of $\beta$.  The order is
$\exists C_2\,\forall m\,\forall^\infty N\,\forall\beta$,
exactly as in \texttt{dualN}, on our selected $a=n$, $g=0$ data.
No inversion of $n$, $n!$, or $3$ occurs.
% END CFO7

% BEGIN CFO8
\paragraph{CFO8. A uniform finite-matrix kernel bound over the DVR.}
\label{claim:CFO8}
Let $f:A^a\to A^b$ be a matrix map of rank $a-r$ over $E$.
Then a constant $c_f\ge0$, independent of $m$, satisfies
\[
 |\ker(f\bmod t^m)|\le Q^{mr+c_f}\qquad(m\ge0).                 \tag{CFO.16}
\]
Choose a nonzero $(a-r)$-square minor $\Delta$ and put
$e=\operatorname{ord}_t(\Delta)$.  For each choice of the
remaining $r$ coordinates modulo $t^m$, differences between
two solutions in the chosen coordinates are killed by the
square minor matrix.  Its adjugate shows that each coordinate
of the difference is killed by $\Delta$.  There are
$Q^{\min(e,m)}$ choices for such a coordinate, so the fiber
has at most $Q^{(a-r)e}$ elements.  Take
$c_f=(a-r)e$.  A zero-rank map has exactly $Q^{ma}$ kernel
elements and is covered by $c_f=0$.

The same bound applies to simultaneous equations indexed by
an arbitrary set if their common kernel over $E$ has dimension
$r$: descending intersections of subspaces of the finite
dimensional $E^a$ reach that kernel after finitely many
equations.  Use this finite subsystem and (CFO.16).  All
choices, including the minor, are made on the integral
equations before $m$ is chosen.  This proof needs no integral
diagonalization and allows the minor to have positive
$t$-valuation.
% END CFO8

% BEGIN CFO9
\paragraph{CFO9. The three invariant bounds with the original old subgroup.}
\label{claim:CFO9}
Put $\nu=\binom n2$.  There exist constants $C_1,C_0,C_j$
depending only on the fixed original data such that, for every
$m$,
\[
 \begin{split}
 |\{Y\in M_m:Y^t=-Y,\ r_gY=Yr_g\ (g\in G_0)\}|&\le Q^{C_1},\\
 |\{Y\in M'_m:\chi(g)r_gY=Yr_g\ (g\in D_0)\}|&\le Q^{C_0},\\
 |\{Y\in\operatorname{Mat}_n(A_m):r_{s_j}Y=Yr_{s_j}\}|
       &\le Q^{mn+C_j}. 
 \end{split}                                                     \tag{CFO.17}
\]
For the first bound, the integral skew matrices form $A^\nu$
because $2$ is a unit.  Over $E$, the matrices $r_g$, for
$g$ in the full cyclotomic kernel, span
$\operatorname{Mat}_n(E)$ by the original full-block
\texttt{NP.span}.  A commuting matrix is scalar: commuting
with $E_{ii}$ makes it diagonal and commuting with $E_{ij}$
equates its diagonal entries.  A scalar skew matrix is zero
in odd characteristic.  CFO8 applies with $r=0$.

For the second bound, the common integral kernel is zero by
the literal $x.N4(D_0)$.  Any vector in its fraction-field
kernel can have its denominators cleared, so that kernel is
zero too.  Apply CFO8 with $r=0$ to the equations
$\chi(g)r_gY-Yr_g=0$.  This is where the old $D_0$ is essential.
It has not been replaced by the bottom subgroup in the
intermediate frame.

For the last bound, the original \texttt{Eig} supplies $n$
distinct roots of the characteristic polynomial in $A$.
Over $E$, eigenvectors for distinct roots are independent
(apply successive factors of the characteristic polynomial
to a proposed relation); hence $r_{s_j}$ is diagonalizable
with distinct diagonal entries.  Its centralizer then
consists of diagonal matrices in that basis and has dimension
$n$.  CFO8 applies with $r=n$ to its commutator map.  The
eigenvector basis is used only over $E$, so no false
integrality assertion about that basis is required.
% END CFO9

% BEGIN CFO10
\paragraph{CFO10. The literal primal and dual actions and a perfect pairing.}
\label{claim:CFO10}
On $\Gamma_0$, and by inflation on $\Gamma_N$, the actions are
\[
 \begin{array}{c|cc}
       &g\in G_0&c\\ \hline
 M_m   &Y\mapsto r_gYr_g^{-1}&Y\mapsto-Y^t\\
 M'_m  &Y\mapsto\chi(g)r_gYr_g^{-1}
                         &Y\mapsto-\chi(c)Y^t .
 \end{array}                                                     \tag{CFO.18}
\]
They are representations.  Indeed $\Gamma_0=G_0\rtimes
\langle c\rangle$ because $\eta(c)=-1$ and $c^2=1$.
The original form says
$r_g^t r_{\sigma g}=\chi(g)^{1-n}I$.
Thus conjugating the adjoint action by negative transpose
gives the action of $\sigma g$; the scalar in this form
cancels.  The dual action satisfies the same relation since
$\chi(\sigma g)=\chi(g)$ and $\chi(c)^2=1$.
Local constancy follows from the finite reductions of $r$
and $\chi$ on the two open cosets.  These are exactly the
two source \texttt{IsAd} actions.  Any representation
satisfying those formulas has the same action.

For $m\ge1$, choose a nonzero $\mathbb F_\ell$-linear
$\lambda:k\to\mathbb F_\ell$ and put
\[
 b_m(X,Y)=\lambda\bigl([t^{m-1}]\operatorname{tr}(XY)\bigr).
                                                                    \tag{CFO.19}
\]
This is a perfect additive pairing.  If a matrix entry has
lowest nonzero term $a t^u$, multiply its transposed matrix
unit by $b t^{m-1-u}$, with $\lambda(ab)\ne0$; such $b$
exists because multiplication by $a\ne0$ is bijective.
This proves nondegeneracy on each side.  Both groups have
cardinality $Q^{mn^2}$, so it identifies one with the full
$\mathbb F_\ell$ dual of the other.  Trace cyclicity and
(CFO.18) give
\[
 b_m(\gamma X,\gamma Y)=\overline{\epsilon(\gamma)}\,b_m(X,Y)
 \qquad(\gamma\in\Gamma_N).                                     \tag{CFO.20}
\]
For $c$, the two minus signs cancel and
$\operatorname{tr}(X^tY^t)=\operatorname{tr}(XY)$.
Thus $M'_m$ is the actual cyclotomic Tate dual required by
\texttt{IsDualPair}; there is no inverse-character error.
At $m=0$ both modules are zero and the unique pairing is
perfect.  This construction uses $\mathbb F_\ell$ coefficients
($r=1$ in the cohomology interfaces), even though the
$t$-length $m$ is arbitrary.
% END CFO10

% BEGIN CFO11
\paragraph{CFO11. The actual decomposition group modulo inertia.}
\label{claim:CFO11}
For a number field $K$, a nonarchimedean prime $v$, and the
chosen extension of its valuation to an algebraic closure
$\Omega$, let $D$ and $I$ be the original decomposition and
inertia subgroups.  Then
\[
 D/I\simeq\operatorname{Gal}(\overline{\kappa(v)}/\kappa(v))
       \simeq\widehat{\mathbb Z}.                               \tag{CFO.21}
\]
The arithmetic Frobenius $x\mapsto x^{|\kappa(v)|}$ corresponds
to $1$.

Here the sole new published import is the finite decomposition
and inertia theorem, in the precise scope checked in the
certificate \texttt{CFO\_LOCAL\_DECOMP}: Milne,
\emph{Algebraic Number Theory}, version 3.08,
Theorem 7.58(a), Proposition 8.10 and Proposition 8.13(b),
\url{https://www.jmilne.org/math/CourseNotes/ANT.pdf}.
For a finite Galois extension of number fields, its
decomposition group is the Galois group of the completions.
Its quotient by inertia is the residue-field Galois group,
and restriction to a Galois subextension maps decomposition
groups onto decomposition groups.  The hypotheses hold here:
we use finite Galois subextensions of $\Omega/K$, a finite
prime, complete discretely valued completions, and finite,
hence perfect and separable, residue fields.

We give the infinite passage explicitly.  The chosen valuation
restricts to a compatible prime on every finite Galois
subextension.  The finite restriction maps just cited are
surjective, and their inverse limit is the subgroup of
$\operatorname{Gal}(\Omega/K)$ preserving that valuation,
namely $D$.  The residue field of the valuation on $\Omega$
is an algebraic closure of the finite $\kappa(v)$.
Indeed each residue element is in the residue field of a
finite extension.  Conversely lift an irreducible polynomial
of any prescribed degree over $\kappa(v)$ to a monic
polynomial over $\mathcal O_K$ and take a root in $\Omega$.
The root is integral at the chosen valuation and its residue
has that prescribed degree.  Finite fields have a unique
extension of each degree inside an algebraic closure.

The finite residue actions are surjective, so compactness
makes the inverse-limit residue action of $D$ surjective;
its kernel, by definition, is $I$.  It induces the first
topological isomorphism in (CFO.21).  Finite-field Galois
groups are cyclic, generated by the stated Frobenius, and
the inverse limit of $\mathbb Z/a\mathbb Z$ is
$\widehat{\mathbb Z}$.  This proves the second isomorphism.
We have not assumed that an infinite extension has a
complete DVR ring of integers.  The author's published
errata were checked; that issue in the proof of 7.50 is not
used in this inverse-limit argument.
% END CFO11

% BEGIN CFO12
\paragraph{CFO12. Build the original arithmetic presentation at every level.}
\label{claim:CFO12}
Fix $N$, and retain the actual $P.V(N)$, $P.e(N)$, $P.v(N,j)$
and $P.B(N,j)$.  The set of allowed finite primes is exactly
the disjoint set indexed by
\[
 \mathcal W=V_\ell\ \sqcup\ s.St\ \sqcup\
                    \{\ast\}\ \sqcup\ \{1,\ldots,q\},             \tag{CFO.22}
\]
where $\ast$ is $s.v_0$ and $j$ denotes $P.v(N,j)$.
The first two sets and $v_0$ are disjoint in the original
setting.  In particular $v_0\nmid\ell$: its residue
characteristic is $s.p\ne\ell$.  The new primes are distinct
and outside $s.S$ by the literal \texttt{v\_inj} and
\texttt{v\_not}.  The equality \texttt{S\_eq} says this list
is all of $P.V(N).S'$.  Thus it is finite and contains every
prime above $\ell$.

At the old primes use the original valuation $s.A(v)$;
at $j$ use $P.B(N,j)$.  For every such valuation $\mathcal A$,
the setting or $P.\texttt{split}$ says
$D_F(\mathcal A)\subset\operatorname{Gal}(\Omega/L)$.
Consequently the literal inclusion
$D_L(\mathcal A)\to D_F(\mathcal A)$ is a continuous
bijection: both groups consist of the same automorphisms
preserving $\mathcal A$.  Compactness and the Hausdorff
topology make it a topological isomorphism.  This is the
actual \texttt{decSplit} dictionary, with no new splitting
choice.

Similarly, \texttt{e\_inj} and the surjectivity field of
$P.e(N)$ identify
$G_{F,P.V(N).S'}$ topologically with $\Gamma_N$.
The \texttt{e}, \texttt{e\_h}, local $\phi$ and $\phi_u$
fields give the identical local maps into that group.
Cyclotomic characters agree because these homomorphisms
preserve $\epsilon$.  Put
\[
 d_v=e(v/\ell)f(v/\ell)=\operatorname{locDeg}(v)
 \quad(v\in V_\ell).
 \]
These numbers are independent of $N$ and sum to $d$.
One direct verification of the sum is to count
$\mathcal O_F/\ell\mathcal O_F$.  The ring $\mathcal O_F$
is a free $\mathbb Z$-module of rank $d$: an integral
primitive element gives a rank-$d$ order, and the
nondegenerate trace pairing places $\mathcal O_F$ inside
its finite trace-dual lattice.  A submodule of that
free $\mathbb Z$-module is free of the same rank.
Thus the quotient has $\ell^d$ elements.
Unique ideal factorization, the Chinese remainder theorem,
and the successive quotients of the localized DVR at $v$
give instead $\ell^{\sum_v e(v/\ell)f(v/\ell)}$.
Equality proves the required sum.

For each real embedding of the totally real $F$, extend it
to $\Omega\to\mathbb C$ and pull back complex conjugation.
It squares to one, acts by $-1$ on every $\ell$-power root
of unity, and does not fix $L$: otherwise the induced
embedding of the CM field $L$ would be real.
After mapping to $\Gamma_N$, these give exactly $d$
elements $c_i$ with $c_i^2=1$, $\eta_N(c_i)=-1$ and
$\epsilon_N(c_i)=-1$.  The original $c_N$ has the same
cyclotomic value.  The count $d$ is the number of embeddings
of a totally real degree-$d$ field, one per infinite place.

The resulting \texttt{ArData} has
$D_q(\ast)=D_L(s.A(v_0))$ and
$D_q(j)=D_L(P.B(N,j))$.  Its map at $\ast$, followed by
$p_N$, is literally the old $V_0.pr$ on this decomposition
group, by \texttt{e\_h}.  Therefore its image is
$V_0.\texttt{frame}.D0$, exactly.  Also $ph_{Nj}$ belongs
to the image of $D_q(j)$ by $P.\texttt{fr\_dec}$ and
$P.\texttt{ph}$; $ta_{Nj}$ belongs to the same image by
$P.\tau_{\mathrm{gen}}$ and $P.\texttt{ta}$, since inertia
lies in decomposition.  This constructs every
\texttt{ArPres} and the three extra range assertions in
the original \texttt{arPres}, simultaneously as a family
over all $N$.
% END CFO12

% BEGIN CFO13
\paragraph{CFO13. Local and global counting from the retained CohP.}
\label{claim:CFO13}
We state explicitly what is being assumed here.  The full
\texttt{CohP} consists of local finiteness, local Tate
duality, local Euler characteristic, global finiteness,
global Euler characteristic, Greenberg--Wiles, normalized
Poitou--Tate and the stated tame local theorem.  This
paragraph uses the first three, global finiteness and
Greenberg--Wiles; it does not prove or replace that
substantial package.

For any finite additive modules with a perfect
$\mathbb F_\ell$-valued pairing as in \texttt{IsDualPair},
both are killed by $\ell$: pair $\ell x$ with everything
and use nondegeneracy.  Thus they are finite
$\mathbb F_\ell$-vector spaces.  A perfect equivariant
pairing identifies the second module with the cyclotomic
dual of the first.  If the first action is locally
constant, so is the second, by this identification and
the continuous cyclotomic character.  Accordingly the
original cohomology inputs apply with coefficient
exponent $r=1>0$.

At a finite prime, write $D$ for its actual decomposition
group and $h^i(M)=|H^i(D,M)|$.  The cohomology groups are
finite by local finiteness.  The definitions of cocycles
and coboundaries give
$|B^1|=|M|/h^0(M)$ and
$|Z^1|=|B^1|h^1(M)$.  Local duality gives
$h^2(M)=h^0(M')$, and local Euler characteristic gives
$h^1=h^0h^2|M|^{d_v}$ at $v\mid\ell$ and $h^1=h^0h^2$
away from $\ell$.  Hence
\[
 |Z^1(D,M)|=|M|^{\,1+d_v}|(M')^D|,                              \tag{CFO.23}
\]
with $d_v=0$ away from $\ell$.  This proves the complete
\texttt{LocCount} contract, including finiteness, for
all the local groups in CFO12.

For the global formula, let $\mathcal L_w\subset Z^1(D_w,M)$
be any additive subgroup containing $B^1(D_w,M)$, and let
$L_w=\mathcal L_w/B^1(D_w,M)$ be its image in $H^1$.
Let $\mathscr S$ denote the Selmer subgroup of $H^1(\Gamma_N,M)$,
and $\mathscr S^\perp$ the group for the orthogonal local
conditions in the dual module.  Global finiteness and
the finite local groups make all these groups finite.
The exact \texttt{CohP.gw} identity, transported through
the actual maps in CFO12, is
\[
 \begin{split}
 |\mathscr S|\,|(M')^{\Gamma_N}|\,
  \prod_w |M^{D_w}|\,\prod_{i=1}^{d}|M^{c_i}|
 =|\mathscr S^\perp|\,|M^{\Gamma_N}|\prod_w|L_w|.
 \end{split}                                                     \tag{CFO.24}
\]
All its hypotheses have been identified: $F$ is a number
field, the set (CFO.22) is finite and contains every
prime above $\ell$, $\ell\ne2$, the actions are continuous,
the chosen valuations lie above their stated primes,
the cyclotomic character is the original one, and the
$c_i$ are conjugations at precisely its real places.
There is no complex-place factor because $F$ is totally real.
Each $L_w$ is an additive, hence $\mathbb F_\ell$-linear,
local condition.

When $\mathcal L_w=Z^1(D_w,M)$, its orthogonal local
condition is zero, by the perfect $H^1$ local Tate pairing.
A cocycle representing a dual Selmer class therefore
restricts to a coboundary there.  Dropping all other dual
conditions only enlarges its set of representatives.
For completeness, let $\mathscr Z$ be the cocycles in
the primal Selmer condition, and let $\mathscr T$ be
the dual cocycles satisfying only the coboundary
conditions at those $w$ with full primal condition.
Use $|L_w|=|\mathcal L_w||M^{D_w}|/|M|$ in (CFO.24),
as well as $|\mathscr Z|=|\mathscr S||M|/|M^{\Gamma_N}|$
and the corresponding identity for the exact dual
Selmer cocycles.  Since $|M'|=|M|$, all global
invariant factors cancel, giving
\[
 |\mathscr Z|\,|M|^{|\mathcal W|}
       \prod_i|M^{c_i}|
 \le |\mathscr T|\prod_w|\mathcal L_w|.
 \]
Multiplying by $|B^1(\Gamma_N,M')|\le |M'|=|M|$
gives precisely the source \texttt{GWBare} inequality,
with the same local maps and $c_i$.
Thus (CFO.23) and (CFO.24) establish the
\texttt{locV}, \texttt{locU}, \texttt{locQ} and
\texttt{gw} fields of the original \texttt{ArData.Counts}.
% END CFO13

% BEGIN CFO14
\paragraph{CFO14. Exact unramified cardinality and completion of ArSys.}
\label{claim:CFO14}
Let $D$ be one of the old tame decomposition groups and
$I$ its inertia.  In the original \texttt{locOf},
\[
 J=I\cap\bigcap_{i=1}^{n}\ker\bar\chi_i.                          \tag{CFO.25}
\]
The residue field of $O$ is finite of characteristic $\ell$
by \texttt{IsCoef}.  The characters $\bar\chi_i$ are
continuous with open kernels.  Thus $I/J$ is a subgroup
of a finite product of multiplicative groups of finite
fields of characteristic $\ell$, and its order is prime
to $\ell$.

For every finite $\ell$-torsion module $M$ with continuous
$D$-action trivial on $J$, let
$Z_J=\{z\in Z^1(D,M):z|_J=0\}$.  Then
\[
                         |Z_J|=|M|.                            \tag{CFO.26}
\]
To prove this, restriction to $I$ factors through $I/J$.
Every cocycle of the latter finite group is a coboundary:
if $z$ is such a cocycle and $a=|I/J|$, summing
$z(gh)=z(g)+g z(h)$ over $h$ gives
$z(g)=gY-Y$ for $Y=-a^{-1}\sum_h z(h)$.
Conversely every coboundary on $I$ is the restriction of
a global coboundary on $D$, which vanishes on $J$.
The restriction image therefore has
$|M|/|M^I|$ elements.

Its kernel consists of cocycles vanishing on $I$.
Such a cocycle takes its values in $M^I$, as follows by
evaluating it at $ig=g(g^{-1}ig)$; it descends to $D/I$.
By CFO11 this quotient is $\widehat{\mathbb Z}$.
Evaluation at arithmetic Frobenius gives a bijection
$Z^1(\widehat{\mathbb Z},M^I)\simeq M^I$.
For surjectivity, if Frobenius acts with finite order
$a$ on $M^I$, a prescribed value $v$ defines the usual
cocycle on $\mathbb Z$; its value on $a\ell$ is
$\ell(1+\cdots+\phi^{a-1})v=0$, and the action on
$a\ell$ is trivial.  The cocycle therefore factors through
the finite cyclic group of order $a\ell$ and extends
continuously to $\widehat{\mathbb Z}$.  Uniqueness follows
from density of $\mathbb Z$.  The kernel of restriction
has $|M^I|$ elements.  Multiplying the two finite fiber
counts proves (CFO.26), including finiteness.

This proves \texttt{UnrCount} for every old tame place
and every module in its original quantifiers.  Together
with CFO12--CFO13 it constructs the entire original
$\texttt{ArSys}(T,d,D_0)$: all levels have their required
\texttt{Counts}; $d_v$ is common to them; the distinguished
range is the old $D_0$; and $ph_{Nj},ta_{Nj}$ are in the
same added local image.  No local or global duality has
been imported beyond the expressly retained
\texttt{CohP}.  The use of arithmetic Frobenius in this
procyclic count does not alter the inverse-Frobenius
elements specified in $P.\texttt{fr\_frob}$.
% END CFO14

% BEGIN CFO15
\paragraph{CFO15. The two scalar local counts.}
\label{claim:CFO15}
Unpack the original \texttt{LocCFT} once, before choosing $m$:
it supplies integers $d'_v$ with $\sum_v d'_v=d$ and the
original \texttt{LocAb} at each $v\mid\ell$.  For any finite
$\ell$-torsion group $V$, a locally constant additive
homomorphism $D_v\to V$ vanishing on the original $\mu_v$
factors, by \texttt{LocAb.univ}, through its surjection
\[
             D_v\longrightarrow\mathbb Z_\ell^{\,d'_v+1}.
 \]
It is determined by the images of the $d'_v+1$ standard
generators.  Indeed any element of $\mathbb Z_\ell$ differs
from an integer representative modulo $\ell$ by an
$\ell$-multiple, which every such homomorphism kills.
No continuity assumption on the factored homomorphism is
needed for this argument.  For $V=A_m$ it follows that the
set of these homomorphisms is finite and has at most
\[
                         Q^{m(d'_v+1)}                         \tag{CFO.27}
\]
elements.  The factorization uses exactly the original
$\mu_v=\texttt{muG}(\ell,D_v)$, not all inertia or its
prime-to-$\ell$ quotient.

For an adically continuous character
$\theta:D_v\to A^\times$ which is not the local
cyclotomic character $\chi$, there is also a constant
$c_\theta$ such that
\[
 |Z^1(D_v,A_m(\theta))|
       \le Q^{m(1+d_v)+c_\theta}\qquad(m\ge0).                   \tag{CFO.28}
\]
Use (CFO.23) and the perfect scalar version of (CFO.19).
The dual character is $\chi\theta^{-1}$.  Choose $g$ with
$a=\theta(g)-\chi(g)\ne0$ and let
$c_\theta=\operatorname{ord}_t(a)$.  Every dual invariant
$y$ is killed by $a$, because
$\chi(g)\theta(g)^{-1}y=y$ and $\theta(g)$ is a unit.
The kernel of multiplication by $a$ in $A_m$ has
$Q^{\min(c_\theta,m)}$ elements.  Formula (CFO.23) proves
(CFO.28), and proves finiteness as well.

For the required off-diagonal characters
$\theta_{ij}=\psi_{v,i}\psi_{v,j}^{-1}$, $i<j$, the
second original \texttt{NP.N1} clause is exactly
$\theta_{ij}\ne\chi$.  Their adic continuity follows from
the original matrix continuity and the diagonal entries
of $w_v^{-1}r_{\phi_v(g)}w_v$; inversion uses $g^{-1}$.
All constants $c_{v,ij}$ are therefore chosen on the
original ordered characters, before $m$ or $N$.
% END CFO15

% BEGIN CFO16
\paragraph{CFO16. Count the actual upper triangular ordinary cocycles.}
\label{claim:CFO16}
Let $Z_{b,v,m}$ be the source $\kappa.\texttt{Zb}(A_m,v)$:
its elements are local primal cocycles which, after
conjugation by the same integral $w_v$, are upper
triangular and have diagonal zero on $\mu_v$.  There is
a constant $C_v$ such that, for all $m$,
\[
 |Z_{b,v,m}|\le
 Q^{\,m\{n(d'_v+1)+\nu(1+d_v)\}+C_v}.                           \tag{CFO.29}
\]
To prove it, write $B(g)=w_v^{-1}z(g)w_v$ and
$R(g)=w_v^{-1}r_{\phi_v(g)}w_v$.  The matrices $R(g)$ and
$R(g)^{-1}$ are upper triangular with diagonal entries
$\psi_{v,i}(g)$ and their inverses.  Order the upper
triangular coordinates $(i,j)$ by increasing $j-i$,
breaking ties arbitrarily.

On the subgroup where all earlier coordinates vanish,
the $(i,j)$ coordinate of
$R(g)B(h)R(g)^{-1}$ is
$\psi_{v,i}(g)B(h)_{ij}\psi_{v,j}(g)^{-1}$.
Indeed a potentially nonzero summand has
$i\le p\le q\le j$.  Unless $(p,q)=(i,j)$, it has
$q-p<j-i$ and is zero by the earlier-coordinate condition.
Thus the coordinate satisfies the scalar cocycle law
for $\theta_{ij}$.
For $i=j$, this is the additive homomorphism law, and
the value on $\mu_v$ is zero by the definition of $Z_b$.

The diagonal coordinate images are bounded by (CFO.27);
the $\nu$ strictly upper triangular images are bounded
by (CFO.28).  The successive kernels of these coordinate
maps end in zero, since all matrix entries then vanish
as functions.  Inducting backwards through this finite
sequence shows the whole group is finite and its
cardinality is bounded by the product of those image
cardinalities.  Take
$C_v=\sum_{i<j}c_{v,ij}$.  This gives (CFO.29).
It counts the actual cocycles, not independent formal
entries, and retains their specified flags and character
order.
% END CFO16

% BEGIN CFO17
\paragraph{CFO17. Ordinary local conditions and the real-place factor.}
\label{claim:CFO17}
Let $\mathfrak b_{v,m}\subset M_m$ consist of the matrices
upper triangular in the same $w_v$ basis.  It has
$Q^{m(n+\nu)}$ elements.  Its coboundaries lie in
$Z_{b,v,m}$: conjugation by $R(g)$ preserves upper
triangular matrices and leaves their diagonal unchanged,
so the diagonal of a coboundary is zero for every $g$.
Let
$\mathcal L_v=B^1(D_v,M_m)+Z_{b,v,m}$, exactly the
ordinary local cocycle condition.  Put
$K_v=M_m^{D_v}$.  The finite sum formula gives
\[
 \begin{split}
 |\mathcal L_v|
 &=\frac{|B^1(D_v,M_m)|\,|Z_{b,v,m}|}
          {|B^1(D_v,M_m)\cap Z_{b,v,m}|}\\
 &\le
 \frac{(|M_m|/|K_v|)\,|Z_{b,v,m}|}
      {|\mathfrak b_{v,m}|/|\mathfrak b_{v,m}\cap K_v|}
 \le \frac{|M_m|\,|Z_{b,v,m}|}{Q^{m(n+\nu)}} .
 \end{split}                                                     \tag{CFO.30}
\]
Every denominator is positive.  The final inequality
uses the inclusion $\mathfrak b_{v,m}\cap K_v\subset K_v$.
This proves the exact local factor that is removed in
the global count.

For each real-place conjugation $c_i$ from CFO12, put
$h=c_i c_N\in G_N$.  Its projection satisfies
$\sigma(p_Nh)=(p_Nh)^{-1}$ and $\chi(p_Nh)=1$.
The form and the representation law therefore imply
$S=r_{p_Nh}$ is symmetric and invertible, over $A_m$.
The primal action of $c_i$ on $M_m$ is
\[
                         Y\longmapsto-SY^tS^{-1}.
 \]
Multiplication $Z\mapsto SZ$ is a bijection from skew
matrices to its fixed submodule, since $S^t=S$.
Consequently, for every $N,m,i$,
\[
 |M_m^{c_i}|=Q^{m\nu},\qquad
 \prod_{i=1}^{d}|M_m^{c_i}|=Q^{m\nu d}.                          \tag{CFO.31}
\]
This is an equality on the selected full-block data.
It does not require diagonalizing a symmetric form or
adjoining square roots.
% END CFO17

% BEGIN CFO18
\paragraph{CFO18. Identify the literal Sel1 and Bd1 with Selmer classes.}
\label{claim:CFO18}
At level $N$, choose the following actual local cocycle
conditions for the primal module $M_m$:
\[
 \begin{array}{c|l}
 w&\mathcal L_w\\ \hline
 v\in V_\ell&B^1(D_v,M_m)+Z_{b,v,m}\\
 u\in s.St&\{z\in Z^1(D_u,M_m):z|_{J_u}=0\}\\
 \ast,\ j&Z^1(D_w,M_m).
 \end{array}                                                     \tag{CFO.32}
\]
All contain the corresponding coboundaries.  At the
tame places the action is trivial on $J_u$: the original
\texttt{NP.ine} kills its image under $\phi_u$, and
$p_N$ preserves that local map.  Thus (CFO.26) applies.
Let $\mathscr S_m(N)$ be the resulting Selmer subgroup
of $H^1(\Gamma_N,M_m)$.
Then there is a natural bijection
\[
 \operatorname{Sel1}_N(A_m,\kappa)/
       \operatorname{Bd1}_N(A_m,\kappa)
                       \ \simeq\ \mathscr S_m(N).               \tag{CFO.33}
\]
We verify both the quotient and every local condition.

A cocycle $y$ on $\Gamma_N$ can be made zero at $c_N$:
subtract $\delta(-y(c_N)/2)$, because
$c_Ny(c_N)=-y(c_N)$.  Its local classes remain unchanged.
A normalized cocycle is determined by its restriction
$z$ to $G_N$, and that restriction has
$z(\sigma g)=-z(g)^t$.
Conversely any primal cocycle on $G_N$ with this
polarization extends by $y(gc_N^e)=z(g)$.
The semidirect-product multiplication verifies the
cocycle law, and the two cosets are open, so local
constancy is preserved.

Since $t^m=0$ in $A_m$, the original entrywise
$t$-adic continuity in \texttt{SelZ} is equivalent to
local constancy; at the $m$th quotient the ideal is zero.
The condition at $u$ is exactly $z=0$ on the image of
$J_u$, which is the literal \texttt{Glob.frame.Iu}.
At $v$, membership in $B^1+Z_b$ says that there is a
matrix $Y_v\in M_m$ for which
$w_v^{-1}(z-\delta Y_v)w_v$ is upper triangular and
has diagonal zero on $\mu_v$.  These are exactly the
two conjuncts of the source \texttt{LocAt}; the same
$w_v$ and $Y_v$ are used.  There are no further
conditions at $\ast$ or the added primes in \texttt{Sel1}.

Finally, a global coboundary $\delta Y$ is normalized
precisely when $c_NY=Y$, namely $Y^t=-Y$.
Its restriction is exactly the source \texttt{Bd1}.
Thus two normalized representatives define the same
Selmer class exactly when their restrictions differ
by \texttt{Bd1}, proving (CFO.33).
Global finiteness from CFO13 makes the global cocycles
finite; alternatively $|Z^1|=|B^1||H^1|$ with both
factors finite.  Hence \texttt{Sel1} and \texttt{Bd1}
are finite as asserted, and
\[
 |\operatorname{Sel1}_N|
     =|\operatorname{Bd1}_N|\,|\mathscr S_m(N)|.                 \tag{CFO.34}
\]
The quotient here is a quotient of additive groups.
No assertion that its dimension vanishes follows from
$\kappa.g=0$.
% END CFO18

% BEGIN CFO19
\paragraph{CFO19. The one-level local/global inequality.}
\label{claim:CFO19}
Fix $m$ and a level $N$ in the eventual set of CFO6--CFO7.
The dual Selmer group for (CFO.32) satisfies
\[
                         |\mathscr S_m^\perp(N)|\le Q^{C_2}.
                                                                    \tag{CFO.35}
\]
Indeed every dual Selmer cocycle restricts to a
coboundary at each added prime, because its primal
local condition there is all $H^1$.  In the same local
group lie both $ph_{Nj}$ and $ta_{Nj}$, by CFO12.
One local coboundary matrix gives both equalities in
CFO7.  The set of such dual cocycles is therefore
contained in $\mathcal S_\beta$, and contains all
global coboundaries.  Dividing (CFO.14) by their
positive cardinality proves (CFO.35).  The other
orthogonal local conditions can only decrease it.

The primal global invariant group has at most $Q^{C_1}$
elements by CFO9: $p_N$ is surjective, so invariance
under $\Gamma_N$ is the same as under $\Gamma_0$.
At $\ast$, the dual local invariants have at most
$Q^{C_0}$ elements by the exact old-$D_0$ image in
CFO12 and CFO9.  At the added prime $j$, invariance
under the local group implies invariance under
$ph_{Nj}$.  By (CFO.13) this is the centralizer
condition for $r_{s_j}$, and so their number is at
most $Q^{mn+C_j}$.  Define
\[
 e_v(m)=m\{n(d'_v+1)+\nu(1+d_v)\}+C_v,\quad
 f_\ast(m)=C_0,\quad f_j(m)=mn+C_j.                             \tag{CFO.36}
\]

Insert (CFO.35) into the exact identity (CFO.24),
cancel the positive factors $|M_m^{D_w}|$ using
$|L_w|=|\mathcal L_w||M_m^{D_w}|/|M_m|$, and use
$|(M'_m)^{\Gamma_N}|\ge1$.  This yields
\[
 |\mathscr S_m(N)|\,\prod_i|M_m^{c_i}|
 \le Q^{C_1+C_2}
       \prod_{w\in\mathcal W}\frac{|\mathcal L_w|}{|M_m|}.
                                                                    \tag{CFO.37}
\]
Each factor on the right is now explicit.
At $v\mid\ell$, (CFO.29)--(CFO.30) bound it by
$Q^{e_v(m)-m(n+\nu)}$.
At an old tame place, (CFO.26) makes it exactly $1$.
At $\ast$ or $j$, (CFO.23) away from $\ell$ bounds it
by $Q^{f_w(m)}$.  Use (CFO.31), then (CFO.34), to get
the integer inequality
\[
 \begin{split}
 |\operatorname{Sel1}_N|\,
 Q^{\,m\{(n+\nu)|V_\ell|+\nu d\}}
 \le |\operatorname{Bd1}_N|\,
 Q^{\,C_1+C_2+\sum_v e_v(m)+f_\ast(m)+\sum_j f_j(m)} .
 \end{split}                                                     \tag{CFO.38}
\]
All finiteness and positivity facts used in these
cancellations have been proved.  This is the original
one-level bound in the selected matrix coordinates.
% END CFO19

% BEGIN CFO20
\paragraph{CFO20. Exact exponent cancellation and the original selected B1.}
\label{claim:CFO20}
Set, before fixing $m$,
\[
                  C=C_1+C_2+C_0+\sum_v C_v+\sum_j C_j.
 \]
The two degree sums are separate inputs:
$\sum_v d'_v=d$ from the retained \texttt{LocCFT},
and $\sum_v d_v=d$ from CFO12.  We do not assume
$d'_v=d_v$ at each place.  Expanding (CFO.36) gives
\[
 \begin{split}
 C_1+C_2+\sum_v e_v(m)+f_\ast(m)+\sum_j f_j(m)
 =m(nd+nq)+C+m\{(n+\nu)|V_\ell|+\nu d\}.
 \end{split}                                                     \tag{CFO.39}
\]
Cancel the common positive power of $Q$ in (CFO.38).
We obtain exactly (CFO.4).
For $m\ge1$, take the finite intersection of the
eventual sets used in CFO6 and CFO7.  The arithmetic
presentation and its cohomology counts were supplied
for every $N$, and all remaining estimates were
uniform in $m$ and $N$; hence this is the only
level restriction.  For $m=0$, both first-order
sets consist of the zero function and the inequality
holds at every $N$.

In the actual two links, $n=4$ and $\nu=6$.
The local ordinary exponent is
$m\{4(d'_v+1)+6(1+d_v)\}+C_v$.
Summing contributes $m(10d+10|V_\ell|)$, while the
real-place and ordinary-normalization factors remove
$m(6d+10|V_\ell|)$.  The $q$ added centralizers
contribute exactly $4mq$; the old distinguished
place contributes only a constant because of $N4(D_0)$.
The remaining exponent is therefore
\[
                            m(4d+4q)+C.                         \tag{CFO.40}
\]
Every division by a scalar in the proof was by a
unit (notably $2$ and integers prime to $\ell$);
the proof includes $\ell=3$.
This supplies the literal selected
$T.\texttt{B1}(\kappa,4[L^+:\mathbb Q]+4T.q)$.
It is a completed conditional implication on the
original data, rather than an unproved first-order
estimate left as an input.
% END CFO20

% BEGIN CFO21
\paragraph{CFO21. Install the coupled supplier without changing the support frontier.}
\label{claim:CFO21}
Apply CFO20 to the precise point/list/extension selected
in CCH15, with CAP13's realization of that same list.
CCH13 preserves the determinant polynomial law, the
point kernel, the ordered local characters and
$N4$ at the original old subgroup under this extension.
The residual determinant argument and its pro-$\ell$
kernel conclusion used in JSP11 and CCH17 apply to
that same point.  Thus the \texttt{PFromAutOf} consumer
now receives together
\[
 \begin{gathered}
 \text{the original one-point \texttt{Choice}},\qquad
 \kappa.g=0,\qquad
 T.\texttt{B1}(\kappa,4d+4q),\\
 T.\texttt{B2}(\kappa)\text{ with }C=0
       \text{ at every truncation and every level}.
 \end{gathered}                                                   \tag{CFO.41}
\]
No new choice of $\kappa$, of its detector, of the
coefficient extension, or of $T$ is made between these
uses.  Consequently the selected
\texttt{JSP\_FIRST\_ORDER\_BOUND} input is removed in the
$h=0$, \texttt{NiceI} support route.  We have supplied
its actual consumer scope; the universal
\texttt{FirstOrderDual}\allowbreak\texttt{ArithmeticBound} on arbitrary
$a$, $g$ and \texttt{Kap} is not being declared proved.

The ideal is still the original $J$ from the same
unitary ordinary system.  The remaining producer
requirements are unchanged: the original ordinary
ring/module/control and arithmetic-fiber data
(including type, neat level and \texttt{TThetaEx});
the four further Hida endpoints on that data; the
full \texttt{HasAut} on this identical $J$ at every
actual Taylor--Wiles realization; and the original
one-point $h=0$ \texttt{Raw} stable-prime propagation.
The entire stated \texttt{CohP}, \texttt{LocCFT}
and arithmetic setting remain explicit hypotheses.
CAP17 supplies \texttt{IsRes} only when its original
type data have been supplied.  CFO adds no ordinary
automorphic existence theorem.

This installation is made separately for both
original rank-four \texttt{Link} objects.  Their
common residual coefficient bound and their distinct
auxiliary primes are preserved.  In the tower the
complex embedding $\iota$ is fixed first; the
original eventual coefficient threshold is then
chosen for both links together, and the construction
must hold for every later coefficient space.  Our
constant $C$ belongs inside the subsequent fixed
point/list estimate, before $m$ and eventual $N$.
It neither creates a separate threshold for
\texttt{HidaPropsAlt} and \texttt{P1core} nor replaces
their common witness $J$ by two ideals.

The contribution is shared by RH and unrestricted
Dirichlet GRH through their Count-II branch.
The target-specialization audit KA5--KA6 still
removes only Case III for characters of
$2$-power order, including the principal target.
It does not make Counts I, II or V vacuous.
Unrestricted Dirichlet GRH still needs its
additional Case-III producer.  None of those
other arithmetic producers is closed by the
first-order argument above.
% END CFO21
